\documentclass[sigconf]{acmart}

%manuscript,review,anonymous
%
% Page Numbers
%\settopmatter{printfolios=true}

% Shrink copyright stuff
%\settopmatter{printacmref=false}
%\renewcommand\footnotetextcopyrightpermission[1]{}
%\usepackage{silence}

\copyrightyear{2026}
\acmYear{2026}
\setcopyright{cc}
\setcctype{by}
\acmConference[DIS Companion '26]{Designing Interactive Systems Conference}{June 13--17, 2026}{Singapore, Singapore}
\acmBooktitle{Designing Interactive Systems Conference (DIS Companion '26), June 13--17, 2026, Singapore, Singapore}
\acmDOI{10.1145/3802974.3809446}
\acmISBN{979-8-4007-2632-3/2026/06}



% = = = = = Packages = = = = = %

%\input{setup/packages}
\input{setup/macros}

% = = = = = Title = = = = = %

\begin{document}

\title{Security Noir: Critical Designs for Security \& Privacy}

\author{Didem Demirag}
\email{demirag.didem@uqam.ca}
\orcid{0009-0003-9708-991X}
\affiliation{\institution{UQAM} \city{Montreal} \state{QC} \country{Canada}}


\author{Leah Zhang-Kennedy}
\email{lzhangke@uwaterloo.ca}
\orcid{0000-0002-0756-0022}
\affiliation{\institution{University of Waterloo} \city{Stratford} \state{ON} \country{Canada}}

\author{Jeremy Clark}
\email{j.clark@concordia.ca}
\orcid{0000-0002-3533-5965}
\affiliation{\institution{Concordia University} \city{Montreal} \state{QC} \country{Canada}}

\begin{abstract}

We introduce \textit{Security Noir}, a critical design approach for engaging with the often hidden and elusive dimensions of security and privacy. We present ten speculative artefacts, expressed through short vignettes that situate each design within everyday contexts of use. These plausible yet provocative scenarios use exaggeration, satire, and juxtaposition to surface socio-technical tensions. We position text-based vignettes as a deliberate design material that enables speculative artefacts to be rapidly produced, shared, and critically engaged with. This work contributes an initial proof-of-concept corpus of critical design artefacts and outlines future work to formalize this approach into a more explicit and transferable design methodology.

%In this paper, we investigate an alternative method for communicating the somewhat elusive security risks and privacy invasions of technology, the web, social networks, and other aspects of modern life. Specifically, we draw from the \textit{critical design} movement popularized by Dunne and Raby in the 1990s. In this project, we have conceptualized 10 new technological designs and described their use through short vignettes. These designs are imaginary but plausible technologies that illustrate a security or privacy risk through absurdity, satire, or humour.

\end{abstract}

\maketitle


% = = = = = Main Body = = = = = %

% !TEX root = ../main.tex

% = = = = = = = = = = = = = = = = = = = = = = = = = = = = = = = = = = = = = = = = = = = = %

\section{Introductory Remarks}

\input{figures/examples.tex}

%\textblue{Opening sentence.}

% Stems from ChatGPT:
%What if the technologies we [use] are quietly reshaping the boundaries of privacy and control?"
% innovations that... betray us
% When privacy is reduced to a setting and security to a feature...
% darker side of convenience
% risks that are normalized and overlooked

%Security is about more than keeping threats out; it is about questioning who we let in, and why. When privacy is a `setting' and security is a `feature,' our research community should continue looking at new ways of promoting insight and awareness over the shifting boundaries of privacy and control. 


%\textcolor{blue}{(LEAH: MAYBE WE CAN OPEN WITH ONE OF OUR VIGNETTES AS A HOOK LIKE THIS? I ALSO ADDED A BIT MORE CONTEXT FOR THE INTRO)}

Consider the following vignette from a collection of speculative artefacts:
%
%\begin{quote}
%\textbf{Eleventh Finger.} Biometric authentication based on fingerprints is generally user-friendly and fast. Eleventh finger is a 3D printed rubber finger, customized with the user’s fingerprint. You can put it on a keychain and use it on cold winter days when you do not want to remove your gloves. Alice lets her friend Bob borrow it when he stays at her house. It can also serve as a backup if the worst happens.
%\end{quote}

\design{Eleventh Finger}{Biometric authentication based on fingerprints is generally user-friendly and fast. Eleventh finger is a 3D printed rubber finger, customized with the user's fingerprint. You can put it on a keychain and use it on cold winter days when you do not want to remove your gloves. Alice lets her friend Bob borrow it when he stays at her house. It can also serve as a backup if the worst happens.}

At face value, the design is a convenience-enabling accessory for biometric authentication; yet, the vignette also challenges common assumptions about fingerprints as secure, embodied identifiers. What happens when bodily identity is lent, misplaced, duplicated, or stolen? The scenario does not present a technical solution; instead, it invites reflection. 

This mode of engagement is characteristic of critical design, which involves speculative artifacts and narratives that do not propose deployable products, but rather provoke reflection and discussion. Early work on critical design is commonly attributed to Dunne and Raby~\cite{DuRa01,Dun05}, and it can also be seen as a refinement of earlier movements in arts and design, beginning with the Italian radical design movement and extending through critical practice in HCI~\cite{Mal17}. 

Designing security software and privacy-enhancing technologies requires anticipating threats that have not yet been exploited, a process in computer security known as \textit{threat modelling} \cite{shostack2014threat}. While valuable, this work is often time-consuming for security experts and difficult for practitioners to learn \cite{shostack_fast_cheap_good}. Prior work has therefore explored bridging speculative design with threat identification, highlighting that threat modelling is inherently speculative, requiring the imagination of possible futures \cite{merrill2020security}. From this perspective, tools from speculative design and design fiction can support early-stage and exploratory threat identification.

These approaches are also more lightweight than traditional frameworks, making them accessible to practitioners with limited security expertise. This aligns with calls within industry and practice for ``fast, cheap, and good'' approaches to threat modelling. For example, Shostack + Associates \cite{shostack_fast_cheap_good} advocate starting with a simple but generative question, “What can go wrong?” Speculative methods operationalize this question through narrative and imagination, helping make abstract or overlooked risks more visible.

%Leah: As reviewers have sugguested, I am replacing these traditional critical design examples with more recent literature on design fictions.
%Consider the example provided in Table~\ref{tab:examples}. In contrast to industrial design, the compass table is not oriented around problem-solving, nor the creation of solutions that are functional, useful, and pleasurable, nor is it meant for mass production. Rather, it offers a subtle criticism by illustrating something---electromagnetic radiation---that typically goes unquestioned. It is designed to invoke a dilemma for the user, raising questions, challenging values, and stimulating debate. Critical design considers these goals to be the function (or `para-function'~\cite{Dun05}) of the concept. %, enabling traditional design approaches to maximizing this functionality. 

%Consider the three examples provided in Table~\ref{tab:examples}. Each describes the design of an object or technology. In contrast to industrial design, these technologies are not oriented around problem-solving, nor the creation of solutions that are functional, useful, and pleasurable, nor are they meant for mass production. Rather they offer a subtle criticism by illustrating something---electromagnetic radiation, the inevitability of death, the ubiquity and stealth of surveillance---that typically goes unquestioned. They are designed to invoke a dilemma for the user, raising questions, challenging values, and stimulating debate. Critical design considers these goals to be the function (or `para-function'~\cite{Dun05}) of the concept. %, enabling traditional design approaches to maximizing this functionality. 

%Critical design has significant overlap with conceptual art, utilizing common methods, and has been disseminated through exhibitions and the art literature. However, critical design works by showcasing things that are plausible, believable, and easy for the reader to situate, requiring design practice. Viewing critical design as art creates challenges. Consider the compass table (from Table~\ref{tab:examples}) in the context of an art gallery exhibit---it is meant to illustrate electromagnetic radiation in everyday life, in one's home, in different positions and orientations within that home---all of this is lost when it sits statically in a gallery. Should critical design be displayed in galleries, a vignette, film, or narrative will often accompany the object to describe its context of use? Security and privacy risks are often intangible and invisible, making concepts difficult to communicate to everyday users. We saw a parallel to electromagnetic radiation in this regard, and considered whether critical design could similarly aid in illuminating cybersecurity risks.

In this work-in-progress, we investigate critical design as a lightweight and narrative-driven approach to communicating elusive security and privacy risks.  This work makes two primary contributions. First, we present a curated corpus of ten speculative design vignettes that surface a range of security and privacy tensions in everyday life. Each artefact is constructed to make a particular socio-technical risk perceptible through exaggeration, satire, or narrative juxtaposition, enabling readers to reflect on a range of speclative security and privacy issues such as the portability of biometric identity, ambient surveillance, and the unintended consequences of personalization. Taken together, the collection demonstrates the breadth of concerns that can be explored through critical design in security and privacy contexts. Second, this corpus serves as a proof-of-concept for our work on developing a structured approach to generating critical design probes for security and privacy. 




%I think this statement can be removed as we now have a different framing
%While the designs are exploratory, their grounding aligns with perspectives from usable privacy and security research that emphasize how users construct understandings of risk through informal channels, including entertainment media~\cite{fulton2019effect} and interpersonal stories~\cite{pfeffer2022replication} play a significant role in shaping users’ awareness and interpretation of privacy and security risks. 





%\paragraph{Methodology.} In order to generate designs, we report on four methodologies. First, we (a) used traditional design methodology to generate 10 designs (Section~\ref{sec:meth1}). We then experimented with the generative AI tool \gpt \gptv. Then we (b) fed our designs to \gpt and asked it to generate new designs (Section~\ref{sec:meth2}). It effortlessly generated designs that initially seemed sound, but under greater scrutiny, were lacking. We then (c) experimented with giving \gpt a sketch of what a possible design could involve and then dialogued with it to develop and refine the design, using it is a collaborative tool (Section~\ref{sec:meth3}). We found this approach produced good results and was less effort than method (a). Finally we (d) educated \gpt on the entire design process (Section~\ref{sec:meth4}), trying to involve it in the earliest stages of design (as opposed to (c) where we involved it after we had an initial sketch of an idea). We found this approach was less effective. In our experience, \gpt is capable of assisting in designs but requires human guidance and intervention. Generative AI's strength is in generating large volumes of ideas, while humans can contribute by curating and refining the best ideas.

\section{Related Work and Research Positioning}

%ADD REFERENCES (these are from the list suggested by reviewers, but I pioriized the more recent ones from DIS, CHI, and other ACM publications)

%Vera Khovanskaya, Eric P.S. Baumer, Dan Cosley, Stephen Voida, and Geri Gay. 2013. "Everybody knows what you're doing": a critical design approach to personal informatics. In Proceedings of the SIGCHI Conference on Human Factors in Computing Systems (CHI '13). Association for Computing Machinery, New York, NY, USA, 3403–3412. https://doi.org/10.1145/2470654.2466467

%Richmond Y. Wong, Jason Caleb Valdez, Ashten Alexander, Ariel Chiang, Olivia Quesada, and James Pierce. 2023. Broadening Privacy and Surveillance: Eliciting Interconnected Values with a Scenarios Workbook on Smart Home Cameras. In Proceedings of the 2023 ACM Designing Interactive Systems Conference (DIS '23). Association for Computing Machinery, New York, NY, USA, 1093–1113. https://doi.org/10.1145/3563657.3596012

%James Pierce. 2019. Lamps, Curtains, Robots: 3 Scenarios for the Future of the Smart Home. In Proceedings of the 2019 Conference on Creativity and Cognition (C&C '19). Association for Computing Machinery, New York, NY, USA, 423–424. https://doi.org/10.1145/3325480.3329181

%Rossi, A., Chatellier, R., Leucci, S., Ducato, R., and Hary, E. (2022) What if data protection embraced foresight and speculative design?, in Lockton, D., Lenzi, S., Hekkert, P., Oak, A., Sádaba, J., Lloyd, P. (eds.), DRS2022: Bilbao, 25 June - 3 July, Bilbao, Spain. https://doi.org/10.21606/drs.2022.681

%@Jeremy and Didem, these examples currently don't match what you have in Table 1 (except for Open Informant): Compass Table, Life Count, and Open Informant? Should they?
Much has been written on critical design, and Malpass provides an excellent survey~\cite{Mal17}. Critical design has long explored speculative artefacts as a means of provoking reflection rather than proposing deployable solutions. For example, Dunne and Gaver's \textit{The Pillow} takes the form factor of a transparent pillow and plays radio signals it picks up from the technology of that time, such as phones, pagers, and baby monitors~\cite{DuGa97}; it challenges users to consider how confidential these signals are. Matsumoto's \textit{Open Informant} is an app that ``searches your communications for National Security Agency trigger words and then sends text fragments containing these words to the badge worn by the user for public display''~\cite{Mal17}. Dunne and Raby's \textit{United Micro-Kingdoms} reimagine the UK as four `supershires' where one is occupied by `digitarians:' extreme technocrats ruled by algorithms in a state of total surveillance.  
%\item \textbf{Alternet.} While not critical design, Gold's Alternet has the similar intention of provoking thought through design~\cite{Gol14}. The internet is re-imagined with users retaining full control over their private data, touching on web tracking and our growing digital dossiers. 

While prior work in security and privacy demonstrates the value of critical and speculative design for exposing hidden assumptions, these approaches have largely been applied within specific domains or participatory contexts. For example, Wong et al.\cite{wong2023broadening} use scenario-based workbooks to elicit user values around smart home surveillance, while Pierce \cite{pierce2019lamps} employs speculative scenarios to examine the social implications of domestic technologies. Khovanskaya et al. \cite{khovanskaya2013everybody} similarly use critical design to interrogate personal informatics systems, surfacing tensions around visibility and self-tracking. In security-focused work, Merrill  \cite{merrill2020security} introduces ``security fictions'' to support reasoning about adversarial futures and threat identification. Rossi et al. \cite{rossi2022if} explore how speculative design and foresight methods can be applied to data protection. These efforts show that speculative and critical design can effectively surface privacy and security concerns through situated narratives, participatory engagement, and domain-specific explorations. In contrast, our work centers on the construction of a curated corpus of standalone critical design vignettes. Rather than embedding artefacts within workshops or elicitation frameworks, we focusing on designing compact, reusable narrative artefacts that independently foreground specific socio-technical tensions.

%While these designs demonstrate the value of critical artefacts in exposing hidden assumptions relating to security and privacy, we are unaware of any previous attempts to build a systematic set of designs central to this topic. However, critical design has been used in other areas of computer science. The uptake of critical design in HCI research is encouraged based on the anticipated value critical design will add to the field~\cite{bardzell_what_2013}. Speculative design is another related concept that HCI researchers rely on. An example is to analyze and question the usage of a certain technology and the consequences of it in the near future~\cite{lawson_problematising_2015}. Ludic design is also utilised to create designs that do not dictate a certain purpose or usage, but rather leave the interpretation to users~\cite{gaver_drift_2004}. Somewhat related to critical design, work on exploring alternative ways (like design fiction~\cite{loureiro-koechlin_vision_2022}, role playing game~\cite{merrill_security_2020}, science fiction and design fiction~\cite{wong_real-fictional_2017}) to explain tricky concepts not only in security but also in other fields like sensing technologies. Studying the metaphors used in CS education gives important insights about how students perceive certain concepts~\cite{harper_conceptual_2024}.

\section{Methodology}

This paper reports the first, exploratory phase of a two-phase process. Our goal in this phase was not to present a finalized reproducible method, but to pilot the development and refinement of a small but meaningful set of vignettes. Rather than relying on purely ad hoc brainstorming, we generated concepts by systematically combining four elements: (1) security and privacy concerns, misconceptions, or principles; (2) domains of application; (3) technological modalities or form factors; and (4) anticipated user tensions or reactions. These combinations were then discussed, interpreted, and iteratively refined into speculative artefacts that aimed to remain plausible while foregrounding critical tension. %The resulting vignette corpus serves as a proof of concept for a more formalized methodology to be developed in future work.

As an example, consider the \textit{Eleventh Finger} vignette. It emerged from a combination of the four elements above. First, we aimed to address a dimension of security beyond privacy. To broaden the range of concerns under consideration, we drew on the STRIDE model~\cite{shostack2014threat} and selected spoofing as the focal issue. Second, authentication processes, including passwords, biometrics, and passkeys, provided an appropriate domain. Third, technologies such as smart locks, smart homes, and 3D printing informed the modality of the design. Fourth, the vignette was developed to foreground a user tension between convenience and loss of control.

In practice, each element was represented on an ideation card created by the authors as sticky notes in an online whiteboarding platform, Freeform. Designs such as \textit{Eleventh Finger} emerged by drawing and recombining cards across categories, then refining the resulting combinations into plausible but provocative vignettes. Not every combination produced a viable design. We typically drew one card from each category and used open-ended discussion to assess whether the combination could support a coherent vignette. Some were too awkward, too literal, or did not surface a meaningful security or privacy tension, and were set aside. Others required reinterpretation, recombination, or simplification before yielding a plausible near-future artefact. We retained concepts that achieved three goals simultaneously: conceptual clarity, plausibility, and distinct critical tension. While this process was effective for exploratory ideation, it still depended substantially on open-ended interpretation and collaborative judgment, and is therefore better understood as a structured pilot than as a finalized reproducible method.

The ten vignettes presented here are those that survived multiple rounds of discussion and revision. A central goal of the second phase of this project is to build on this exploratory process and formalize it into a more explicit and transferable design methodology.


%
%
%The ten vignettes presented here are those that survived multiple rounds of discussion and revision.
%
%%Jeremy and Didem, this part is for the methodology details you'll add. Please check for accuracy and edit to match what you write in that section.
%
%Rather than relying on purely ad hoc ideation, our process was guided by the systematic combination of four elements: (1) security and privacy concerns, misconceptions, or principles; (2) domains of application; (3) technological modalities or form factors; and (4) anticipated user tensions. \textblue{Tracked in freeform.} By recombining these elements and iteratively refining resulting concepts, we produced a diverse set of artefacts that balance plausibility with critical tension. The resulting designs suggest that this combinatorial approach can support the generation of varied and conceptually distinct probes, offering a foundation for a more formalized methodology in future work. 
%
%\textblue{merge this with below}
%\textblue{add one example}
%
%Building on these traditions, this work explores how curated collections of critical design vignettes may function as communicative probes for reflecting on security and privacy risks. We structure this research as a two-phase design inquiry intended to move from exploratory artefact generation toward methodological formalization.  The first phase focuses on creating and iteratively refining an initial corpus of speculative designs.
%% through traditional ideation processes and discussions to examine the feasibility, scope, and expressive range of vignette-based critique, specifically assessing whether artefacts could be generated consistently, address diverse dimensions of digital risk, and communicate scenarios in ways that invite reflection and speculation. 
%The second phase (future work in \S\ref{sec:future}) will develop a reproducible design procedure, articulating explicit steps and documentation practices to support adoption or adaptation by other researchers.
%
%This paper reports on the first exploratory phase. Authors with cybersecurity research backgrounds collaboratively generated candidate artefacts through structured brainstorming informed by domain expertise. Ideation processes were captured on sticky notes on a virtual whiteboard and were organized into categories to scaffold ideation, including (1) security and privacy concerns (\eg `if you don't pay for the product, you are the product'), misconceptions (\eg `nothing to hide'), or principles (\eg `minimal disclosure')  that could be the basis of a critique; (2) domains where users might especially care about security or privacy (\eg health records or social networks); (3) a form factor for the new technology (\eg haptics or extended reality); and (4) a possible user response (\eg amused or embarrassed). After all ideas were captured and written as notes, they were randomly combined.  We then drew notes at random (\textcolor{red}{one note from each category}) and held open-ended discussions about what the notes meant to us. \textcolor{red}{Looking at these different combinations helped us brainstorm unconventional combinations to come up with compelling designs. We also individually brainstormed and shared our findings with the other authors in order to challenge the ideas we had developed.} The combination of concepts was discussed in open-ended sessions to explore interpretive possibilities and identify promising narrative directions.\textcolor{red}{During our open-ended discussions, when we could not come us with a design that meaningfully combined the notes we have drawn, we decided to focus on only one or two notes to ensure the designs met our criteria.} Design concepts were then iteratively developed with the aim of producing near-future artefacts that remained plausible while foregrounding critical tension: \eg showing the darker side of convenience, risks that are normalized and overlooked, or ways a user might be betrayed by the technology. Concepts that failed to achieve conceptual clarity, plausibility, or distinctiveness were discarded. The resulting ten designs represent artefacts that survived multiple rounds of discussion and refinement, each requiring substantial time and revision. \textcolor{red}{Either during the design of the artefact or after the vignette is finalized, we discussed potential titles that both represent the artefact and are engaging.}. Next, we present our designs.

% JC: this is all AI related, I think we can drop it or selectively bring back aspects

%\paragraph{Methodologies for Design Creation and Evaluation} 

%While design process can be based on well-established methods, that are relevant to us, like whiteboarding, using sticky notes, affinity diagrams; the integration of AI to the creation process has been proliferated with the abundance and the variety of AI-based tools. These tools can be used to generate the content itself or help the designer along the process of creation and ideation. There are also works that examine how users prompt generative AI (~\cite{sanchez_examining_2023}, ~\cite{chang_prompt_2023}); utilize generative AI for co-ideation or co-creation~\cite{tholander_design_2023},\cite{chiou_designing_2023}.
%
%There are also efforts towards creation of computational tools to that would assist task that are manually done in general. One line of work is about creating tools to group design aspects to help ideation and exploration of design solutions\cite{huang_designnet_2023}. Speculative design is also used to imagine the future of the tools based on generative AI~\cite{lin_generative_2023}. Custom-made AI tools are proposed for specific design goals like integrating context awareness\cite{fan_contextcam_2024}.
%
%While we only created the designs and did not apply any techniques to evaluate them, there are several methodologies for this purpose. User studies like interviews and surveys are one of the most common techniques employed~\cite{sanchez_examining_2023}, ~\cite{chang_prompt_2023}. They are deployed to evaluate different aspects of the design process like assesing the utility of the AI-tools used during the design process, understanding users' needs or deploying a participatory workshop. Methods to asses the creativity of AI-generated content are also proposed~\cite{chakrabarty_art_2024}.

%\paragraph{Alternative Pedagogy} 





%CS education literature has explored how to teach concepts related to computer science in general or cybersecurity to a diverse audience. Some of which are understanding how CS students perceive AI and cybersecurity~\cite{ojha_computing_2023}, using visual cryptography for explaining concepts to K-12 students~\cite{rayavaram_designing_2023}, how to initiate a security mindset~\cite{holley_teaching_2023}. Studying the metaphors used in CS education gives important insights about how students perceive certain concepts~\cite{harper_conceptual_2024}.



% = = = = = = = = = = = = = = = = = = = = = = = = = = = = = = = = = = = = = = = = = = = = %
% Method A

%\subsection{Process}

%\begin{figure*}[t]
%\begin{center}
%\includegraphics[width=0.8\textwidth]{figures/brainstorming.pdf}
%\caption{Final affinity board after brainstorming. Not intended to be legible; intended to communicate the scale of the board.\label{fig:board}}
%\end{center}
%\end{figure*}




\section{Initial Set of Designs}
\label{sec:meth1}



%\textcolor{red}{New methodology}
%\begin{itemize}
%	\item We started by analyzing the different methods of design (divergent, convergent, mapping, prototyping...)
%	\item We decided to follow a divergent method:  the \textit{How might we?} (HMW) technique from design thinking. starting from the general problem statement, each of has created \textit{How might we?} cards. The cards reflected different aims we had for the design we will create like how might we provoke reflection of societal and ethical issues, make the narratives more relatable or create ideation cards that can eventually helps wit coming up with the vignettes. We also tried ChatGPT for creating the card (\textcolor{red}{Any reflections on how it went?}). We then grouped the HMW cards around similar topics and applied dot voting to determine the most important questions.
%	\item We also gathered images, news, concepts etc that reminded us of the project and inspired us in an inspiration board.
%	\item We grouped the HMW cards we created (with and without help of ChatGPT) under themes. We came up with five suits (themes) for ideation cards (Critical, design, scenario, narrative, stickiness). Then we decided to keep only critical, which determines the main topic to be critiqued, and design, which determines the medium. Then redistributed the cards under these two themes. The rest of the themes serves more as a checklist for a quality check for the designs created (for evaluating the narrative style, whether the vignette is memorable and powerful).
%	\item We divided the themes among ourselves so that each of us had cards from each suit (critical and design). Then we each created ideation cards based on the theme with the format: title and description.
%	\item We then tested the cards. Each of us drew one card from each suit and came up with a design. We repeated this to see how we draw cards (e.g., keep drawing 2 cards until we find an idea, keep the one we like and continue drawing from the other suit, draw more than 2 cards from each category) and whether the combinations lead to ideas.
%	\item We then refined the cards. Each of us looked at all of the cards to either present alternatives to the card content itself (changes regarding the title and/or the description) or we suggested to keep them as is. We also merged/unmerged cards, discussed whether the card is a good fit for its suit. 
%	\item After the initial refinement, we applied the MoSCoW (Must Have, Should Have, Could Have and Won't Have this time) method to determine the most important, relevant cards to keep in the deck.
%	\item After several iterations of refinement of the cards involving merging, deleting, moving the cards between the two categories and rewording; we finalized a deck of 37 critique and 30 design cards.
%	\item We conducted an autoethnography for our design process (\textcolor{red}{See section XX}). Each of us created a design a day using the ideation cards. We created a template for our diary entries. We each created a first set of diary entries, and then we refined the fields in each entry to provide a comprehensive overview our design process.(\textcolor{red}{Add the fields we filled each time, also add the total number of design, how long it took etc.})
%	\item We also tried co-designing and creating three designs each, using the same two cards (one from design, one from critical)
%	\item We performed thematic analysis on our diary entriess (\textcolor{red}{See section XX}).
%\end{itemize}
%\subsection{Designs}

\design{No(i)sy elevator}{Elevator music has never been regarded as elevating your mood. This is why employees will be happier riding the noisy elevator. When it recognizes a user, it spins up a selection from their most-played songs on major music streaming platforms. Alice finds hearing her favourite Miles Davis song grounds her at the start of each workday, while Bob and Carol, coincidentally riding the elevator together, discover their shared love for 90s Britpop. David is a bit more concerned about the profanity-ridden banger that the nosy elevator plays for him and his boss.}

\design{Wifi Projector}{Minimalist wifi routers have a single light indicator to show internet connectivity, such as a green light as opposed to red light. This understates the interesting and intriguing data flowing through the device as users browse the internet. By contrast, the wifi projector is a talkative router that projects visualizations onto the ceiling above where it is placed. Every website that pulls in scripts and cookies from other domains is displayed as a constellation of stars of various colours and sizes, which slowly fade away. While Alice's visit to dictionary.com is a beautiful galaxy, she is also perturbed by the number of ads and tracking companies surveilling her.}

\design{Crystal Avatars}{Crystal avatars are created instantly but only get better over time. At first, they are not recognizable and cannot be differentiated from other users, but as users browse the site, small details crystallize. Alice is happy to save time registering for the site and have a personal avatar, which automatically adds a cat accessory after she spent a lunch hour reading about hairball treatments. However, she grows uneasy about the "Woman, Life, Freedom" badge, a cause she believes in but only vocalizes discriminately.}

\design{Chatty}{The latest software update for Alice's voice-activated home assistant Chatty adds machine learning to better infer Alice's commands, such as adding items to her shopping list and suggesting music she might like. It does not always know whether Alice is talking to it or not, so over time it picks up on pieces of Alice's life and has trouble unlearning them. Alice is startled but pleased when Chatty chimes in with the name of `that actor who was in that thing' she was telling her friend about. She is bemused to find a specialized vinyl cleaning solution on her shopping list after she said Bob's records smell fishy.}

\design{Dynamic Laughtrack}{Digital television content, such as sitcoms, encode laughtracks as a series of cues rather than an actual recording. Smart TVs listen and classify the viewer's level of laughter on a 10-point scale (with fine-grained training over time). The TV adds laughter to the content in a sensitive and considerate manner, where the laughtrack is only marginally higher on the scale than the viewer's own laughter. This gently nudges the viewer to greater enjoyment of the program without bombarding her. It also mixes in actual past recordings of the viewer's own laughter to capitalize on emotional mirroring.}

\design{Receiver plant}{The receiver plant in the front yard of Alice records the information of the people that pass by. It is a greedy plant that needs to absorb enough bluetooth data to grow. Bob is singlehandedly responsible for its lustrous canopy from walking by Alice's house every day, mostly while checking his work emails on his phone. Perhaps the plant overstepped when it displayed, "you are 15 minutes late today, are you sure you can make it to work on time?"}

\design{GOTTCHA}{GOTTCHA is a human-detection system to protect online accounts being accessed by bots. It invokes the device camera to analyze if a live human is using the device. To protect against video replays or machine-generated video, it unexpectedly prompts the user with a randomly selected image and carefully measures their reaction to it. Micro-expressions are involuntary facial displays of emotion that are too fast to mimic artificially. GOTTCHA's image bank can provoke disgust, anger, fear, sadness, happiness, surprise or contempt. Over time, Alice stops visiting websites that use GOTTCHA because of its capricious tendency to mix in disturbing images.}

\design{ToS Fishing}{ToS Fishing is a digital game where players earn points by hunting down excessive legalese on websites, including terms of service (ToS), privacy policies, and cookie policies. To play, users simply copy and paste the URL of the legalese. If the ToS has been seen before, the users are awarded points based on its word count. If it has not been seen before, it is reviewed by a human for validity and word count—and in this case, the user gets a finder's bonus. To mitigate people from launching custom websites, the game only accepts sites from the Alexis top million. Users display their aggregate score on a leaderboard with a profile showing their top catches (word count only, not the website), and can earn various badges for playing consistently.}

\design{Eleventh Finger}{Biometric authentication based on fingerprints is generally user-friendly and fast. Eleventh finger is a 3D printed rubber finger, customized with the user's fingerprint. You can put it on a keychain and use it on cold winter days when you do not want to remove your gloves. Alice lets her friend Bob borrow it when he stays at her house. It can also serve as a backup if the worst happens.}

\design{Calm Watch}{Calm Watch is a smartwatch that tracks anxiety levels based on several biomarkers. When the anxiety level of the wearer crosses a threshold, a haptic vibration pulse prompts the user to look at the watch. From there, they can select a variety of calming techniques, including breathing exercises and guided meditation. Bob is nervous to talk to Alice and becomes flustered when his watch starts vibrating mid-conversation, just loud enough for Alice to hear it and notice his fidgeting with the watch.}


\subsection{Discussion}

%REFERENCES TO ADD:
%Lisa P. Nathan, Batya Friedman, Predrag Klasnja, Shaun K. Kane, and Jessica K. Miller. 2008. Envisioning systemic effects on persons and society throughout interactive system design. In Proceedings of the 7th ACM conference on Designing interactive systems (DIS '08). Association for Computing Machinery, New York, NY, USA, 1–10.

\paragraph{Vignettes.} Critical designs are often realized as physical or visual artefacts that are built, exhibited, or otherwise materialized to support engagement and interpretation. For example, Dunne and Raby’s ``Compass Table'' in Table~\ref{tab:examples} (and the seven other objects in their \textit{Placebo} project~\cite{DuRa01}) were fabricated and placed into participants’ homes for everyday use, followed by exit surveys capturing how individuals interpreted and experienced the objects. At the same time, critical designs are frequently accompanied by narrative descriptions that situate the artefact within a context of use, making their intended critique legible.

In this work, we adopt text-based vignettes not as a substitute for design, but as a deliberate design material. Each vignette describes both the form and situated use of a speculative artefact, enabling readers to imagine its integration into everyday life and engage with it critically. When successful, these descriptions support what Malpass describes as the `rhetorical use' of critical design~\cite{Mal17}, where the artefact functions to provoke reflection rather than enable practical deployment. The designs are intentionally transgressive, crafted to evoke emotional and interpretive responses through humour, satire, ridicule, poetry, playfulness, lewdness, appropriation, irreverence, or deliberate overcompensation.

Our use of text aligns with traditions in design fiction and scenario-based design, where narrative is treated as a legitimate medium for constructing and encountering speculative artefacts. Design fiction, in particular, often employs “diegetic artifacts”; objects situated within a fictional world that make speculative technologies tangible and support reasoning about their implications, whether or not they are fully realized as physical artefacts. Similarly, value scenarios \cite{nathan2008envisioning} use textual narratives to explore the broader societal and ethical implications of design decisions. In this sense, our vignettes can be understood as lightweight, narrative instantiations of critical designs that foreground use, context, and consequence.

Vignettes are also widely used in research contexts, defined as ``short stories about hypothetical characters in specified circumstances, to whose situation the interviewee is invited to respond''~\cite{finch1987vignette}. In HCI and usable security research, such narrative constructions are used to surface assumptions, value tensions, and reflections on possible futures when direct observation or deployment is infeasible, unsafe, or premature (e.g.,~\cite{naeini2017privacy, hadan2024privacy}). By leveraging text as a design material, our approach prioritizes accessibility, scalability, and interpretive flexibility, allowing speculative artefacts to be rapidly generated, shared, and adapted while still supporting meaningful critical engagement.

%\paragraph{Vignettes.} Critical designs are commonly described through vignettes, though they are occasionally built and displayed as physical artifacts. For example, Dunne and Raby's ``Compass Table'' in Table~\ref{tab:examples} (and the seven other objects in their \textit{Placebo} project~\cite{DuRa01}), were built and given to individuals to use, followed by exit surveys capturing how the people felt about the objects. We presented our designs as short vignettes that describe both their form and use context. When successful, such descriptions allow readers to readily imagine these artefacts within their own lives and engage with them critically, resulting in their `rhetorical use'~\cite{Mal17} rather than practical deployment. These designs are intentionally transgressive, crafted to provoke emotional responses by invoking humour, satire, ridicule, poetry, playfulness, lewdness, appropriation, irreverence, or deliberate overcompensation.

%Vignettes used in research contexts, in comparison, are defined as ``short stories about hypothetical characters in specified circumstances, to whose situation the interviewee is invited to respond'' ~\cite{finch1987vignette} are established instruments in speculative design, design fiction, and scenario-based inquiry in HCI. In these contexts, narrative depictions of hypothetical technologies or practices are used to surface assumptions, value tensions, and reflections on possible futures without requiring the artefact itself to exist materially. Design fiction, in particular, is not treated as empirical stimuli for behavioural measurement, but as ``diegetic artifacts'' that enable participants to reason about situated technological futures, their consequences, and social norms. Vignette-based methods in usable security research (e.g., ~\cite{naeini2017privacy}, ~\cite{hadan2024privacy}) similarly construct narratives to investigate user expectations, decision-making processes, and ethical considerations surrounding emerging technologies when direct observation or deployment is infeasible, unsafe, or premature.

\paragraph{Usage Contexts.} In this work, we propose vignettes as a methodological vehicle for communicating speculative artefacts rather than as experimental instruments intended to elicit measurable responses. The vignettes function as shareable analytic artefacts that others may appropriate for reflection, inquiry, and design exploration. This approach therefore aligns with established HCI practices in which scenarios are used to probe socio-technical implications. 
%However, it does not aim to produce generalizable behavioural data or predictive claims about user action. Instead, our objective is interpretive and exploratory---to expose assumptions, provoke critical reasoning, and broaden discourse around possible digital futures.

Beyond their role as speculative artefacts, we foresee several practical use cases of critical design vignettes for security and privacy. First, they may serve as reflective prompts in educational, participatory, or professional contexts, where speculative scenarios facilitate discussion of assumptions, values, and mental models surrounding security and privacy practices. Second, they may function as research instruments in future speculative design or vignette-based studies, where they can be adapted as stimuli to explore perceptions of risk, trust, or technological futures. Finally, we envision their use as generative tools in design processes to inspire ideation, highlight overlooked attack vectors, or prompt reconsideration of design assumptions before technologies are implemented.

We therefore position `Security Noir' (an homage to Dunne and Raby’s term `Design Noir'~\cite{DuRa01}) as a critical design approach for engaging with the often uncomfortable or hidden aspects of security and privacy technologies and their impacts on human experience. Methodologically, this approach takes the form of a lightweight, imaginable, and scalable medium that allows speculative artefacts to be encountered without technological implementation, while remaining grounded in established HCI traditions of using vignettes in survey research, design fiction, and scenario-based studies. These critical designs intend to activate readers’ existing mental models and experiences, supporting reflective engagement with security and privacy risks that are otherwise difficult to visualize or experience directly.

%We will provide some of our \textit{stems}\footnote{In this context, `stems' refer to the core ideas, influences, or conceptual elements that contribute to the development of a design.} for one example: Eleventh Finger.  We felt that many of our designs were privacy-focused and we wanted a design that was explicitly security-focused. STRIDE is a threat modelling framework that outlines six high level categories of threats: Spoofing, Tampering, Repudiation, Information Disclosure, Denial of Service, and Elevation of Privilege~\cite{shostack2014threat}. We considered possible designs for each and Eleventh Finger came from our ideation around Spoofing. We tried to avoid privacy implications, even though they are important as well (the ACLU and EPIC were raising concerns about Disney's use of biometric fingerprint scanning at its parks). Biometrics are often considered the holy grail of authentication by non-experts for their simplicity, however they have drawbacks in the areas of entropy, revocation after theft, and delegation. At the hacker conference Def Con, researchers showed that biometric scanners, like those found in modern smartphones, can be tricked using 3D-printed fingerprints~\cite{levalle2020biometric}. We were open to using (dark) humour to avoid sounding too earnest in our designs. The reader can likely see how Eleventh Finger can emerge from these stems, however it is less apparent where the stems themselves came from (\eg from our academic backgrounds and personal experiences), nor the large number of stems we discarded because we could not see connections between them. 


% = = = = = = = = = = = = = = = = = = = = = = = = = = = = = = = = = = = = = = = = = = = = %
% Method B
%
%\section{Methodology B: LLM Generated Designs}
%\label{sec:meth2}
%
%\subsection{Methodology}
%
%The design concepts in the previous section were developed before or independent from generative AI tools like \gpt. However, once these tools became available, we felt it would be interesting to determine how they might assist in generating designs. In particular, since we are developing written vignettes, large language models (LLMs) are the appropriate machine learning technique suited to this purpose. We began with the simplest methodology: we fed \gpt \gptv example designs from Section~\ref{sec:meth1} and asked for the generation of new designs. All transcripts are available on our GitHub repository.\footnote{GitHub: \url{https://github.com/MadibaGroup/2017-Critical-Design}}
%
%\subsection{Designs}
%
%The following are the first two designs generated by \gpt.
%
%\design{Transparent Wallet}{Alice's digital wallet automatically categorizes her purchases and shares anonymized spending habits on a public leaderboard. She enjoys seeing her eco-friendly purchases earning her top marks in the green consumer category. However, Alice is troubled when she notices ads for luxury items appearing more frequently after her high spending months are tracked and shared by the wallet~[\href{https://chatgpt.com/share/4be6016c-7eb0-4063-936a-2a145e925a19}{\gpt Transcript}].}
%
%\design{Echo Chamber Filter}{A social media app uses an advanced algorithm to filter content to match users' perceived political views and biases. At first, Alice appreciates the streamlined news feed with content she agrees with. Over time, she realizes her interactions are limited to a narrow worldview, reinforcing her existing beliefs and making it difficult to encounter differing perspectives or engage in meaningful discussions~[\href{https://chatgpt.com/share/4be6016c-7eb0-4063-936a-2a145e925a19}{\gpt Transcript}].}
%
%\subsection{Discussion}
%
%Transparent Wallet can be considered an example of critical design but Echo Chamber Filter is not. The latter makes a valid critique but it does not describe a new (speculative) technology, personalized feeds have been utilized by social media platforms for years and have been thoroughly criticized for exactly this point (\eg~\cite{snakeoil}). We asked for five designs and three of the designs, including Transparent Wallet, were based on the threat of data sharing being weaponized against the user (twice in the form of ads and once in the form of insurance premiums), demonstrating that \gpt can be highly repetitive. 
%
%Of all the five generated designs, Transparent Wallet was the closest to our own designs. It uses a similar structure: the design, plausible example of its intended use, and an unintended misuse. Gamification, data sharing, and algorithmic inference are all stems we ourselves considered (and even used) for our design. However we would work on improving what is being critiqued, which is pedestrian and straight-forward, strengthen the plausibility of Alice's willingness to share obviously sensitive personal data, and try to make the narrative clearer (we read it several times to `get it'). This led us to consider a methodology where we co-designed with \gpt to iterate and refine designs.
%
%
% = = = = = = = = = = = = = = = = = = = = = = = = = = = = = = = = = = = = = = = = = = = = %
% Method C
%
%\section{Methodology C: LLM Collaborative Co-Designs}
%\label{sec:meth3}
%
%\subsection{Methodology}
%
%We began each design session with a new window in \gpt and started each session with the same prompt: a few sentences about using critical design for security and privacy (assuming it could reference for itself what these terms mean) and our manually created designs from Method A as samples. We then input a set of stems we had brainstormed and asked for ideas about possible designs. As an LLM, we did not concern ourselves with making the stems coherent or even grammatically correct---just enough words to trigger data around the subject. An example set of stems (which was used for the design ElephantMind below) are:
%
%\design{Stems}{Avoiding small talk; to know more about the person you meet; the app reminds you the last conversation you had with them; constantly reminding something trivial (Alice promised Bob to grab a coffee sometime. Next time Bob sees Alice the app reminds Bob what Alice said.); Summarize the last conversation you had with them; For a person you meet for the first time, device gives you some background info about them so that you choice the right topic of conversation.; Name for the app: something that is about never forgetting like having a memory like an elephant; Photographic memory; Or something related to recording everything}
%
%We then asked for examples along the lines of: technologies that could serve as the form factor for a critical design, potential names for the design, and examples of `things that could go wrong' with the technology. In most cases, we would ask for 10 examples. We would consider the examples and apply redirection if the examples were not what we were looking for; if they were, we would curate 1 or 2 that were workable and ask for further examples. At certain points, we would ask it to write a draft vignette and then manually edit it.
%
%\subsection{Designs}
%
%\design{BloodKey}{BloodKey is a form of two-factor authentication that can protect online accounts from stolen passwords. With BloodKey, after providing their password, users authenticate with a pin-prick blood test using an inexpensive biomedical Bluetooth gadget. BloodKey was quickly adopted by leading websites because it was very hard to fake, despite its invasive nature. However, Alice grew tired of the daily inconvenience of pricking her finger, so she gradually stopped logging into her social media accounts altogether, while Bob's medical condition prevented him from using it at all. A media frenzy ensued when the most popular free-to-play online game sold users' genomic data to pharmaceutical companies, boosting their revenue streams~[\href{https://chatgpt.com/share/1421429f-6310-4bd7-8e83-202f9e8108ed}{\gpt Transcript}]. }
%
%\design{Civic Compass}{Civic Compass is an add-on for popular GPS navigation apps that highlights the ethical considerations of various routes, such as avoiding surveillance hotspots or environmentally protected habitats. Public policy debates have emerged over the app's potential abuse by criminals. When Alice is in a rush, she ignores its suggestions. She is also disturbed to see that most of the surveillance hotspots are in impoverished, marginalized neighborhoods~[\href{https://chatgpt.com/share/4814bd4f-bfcd-485b-9f19-78889ee93eb9}{\gpt Transcript}].}
%
%\design{ElephantMind}{ElephantMind is an augmented reality (AR) social interaction app designed to enhance small talk by reminding users of past conversations. Alice cancels her coffee date with Bob when ElephantMind's service is down, fearing she won't remember their last chat. When Alice enthusiastically asks John about his trip to Paris, not knowing he was there for a difficult family matter, it leads to an awkward moment. Bob feels their friendship weakening as ElephantMind constantly pushes him to talk about trending political issues with Alice~[\href{https://chatgpt.com/share/237f70de-d08d-4830-b8f8-977de7512807}{\gpt Transcript}].}
%
%\design{Profiled Plates} {Profiled Plates is a technology for high-end restaurants that offers an extraordinarily personalized dining experience by using extensive data mining and online profiling. With just a reservation name, the restaurant gathers detailed information about guests and their friends from social media, public records, and online interactions to create a bespoke meal tailored to individual preferences and backgrounds. Alice books a table for her birthday and is thrilled to find her favorite flowers on the table and a menu that excludes cilantro, a dislike she once tweeted about. However the system led to a low tip from Bob when it served a meal highly customized for a different Bob Bradley~[\href{https://chatgpt.com/share/5a541220-b6d4-4faf-9b2e-4a979a6bbdf1}{\gpt Transcript}].}
%
%\design{Square One}{Square One is a smart deletion feature for social media accounts which resets interests, deletes content, and unfollows people based on a simple prompt from the user about what they want removed. Jack uses Square One to unremember a volunteer project he led, only to find that potential employers now see him as less community-minded and pass him over for jobs. Jack is dismayed as he watches these projects falter under less-qualified hires~[\gpt Transcripts \href{https://chatgpt.com/share/fdf2764b-2c8b-42d5-aada-bfc5f0d0dcac}{A} and \href{https://chatgpt.com/share/ff5c6451-9fb4-4021-88b9-e4df181d6776}{B}].}
%
%\subsection{Discussion}
%
%Our experience with this methodology was positive. \gpt excelled at providing countless examples and as long as it produced a few useful ones, we were able to work with it in refining them. We found this methodology required less effort than traditional design (A) and produced results at least as good. One issue we found is when we used new windows and settings to limit cross-contamination between sessions, we found \gpt seemed to remember past interactions. The primary drawback of this method was needing a set of stems before interacting with \gpt. For this reason, we decided to attempt a final methodology where we involved \gpt from the beginning of the design process. 
%
% = = = = = = = = = = = = = = = = = = = = = = = = = = = = = = = = = = = = = = = = = = = = %
% Method D
%
%\section{Methodology D: LLM Design Methodology}
%\label{sec:meth4}
%
%\subsection{Methodology}
%
%\textblue{To be completed again. We performed this methodology but will revisit it with more time and effort.}
%
%In the previous methodology, we provided \gpt with stems and found success in iterating these stems into designs using \gpt as a collaborator.  This led us to question whether it would be useful to collaborate with \gpt from the very beginning of the design process, taking its input on setting up the design methodology, on finding additional sticky notes or different clusterings, and ideating through different combinations of sticky notes.
%
%\begin{itemize}
%\item Prompts about what methodology to use
%\item Prompts for help in clustering notes in affinity map
%\item Prompts for additional notes in each cluster
%\item Prompts for generating stems from combination of notes
%\item Prompts for expanding stems
%\item Once we have stems, go to Methodology C
%\end{itemize}
%
%\subsection{Designs}
%
%\textblue{To be completed}
%
%\subsection{Discussion}
%
%\textblue{To be completed}
%


\subsection{Future Work}
\label{sec:future}

%Leah: revised this section to address reviewer comments. The user study bit is added (3rd point).
As a work-in-progress, this paper reports on the initial exploratory phase of our approach. Future work will focus on formalizing the design process into a reproducible method, expanding the space of concerns and domains considered, and evaluating how such vignettes function in practice (e.g., in workshops, educational settings, or participatory design contexts). Through this progression, we aim to establish critical design vignettes as a viable methodological tool within usable security and HCI research. 

First, while this work demonstrates our approach based on concerns, domains, modalities, and user tensions, future work will expand and refine these dimensions through a more comprehensive synthesis of usable security and privacy literature. This includes developing more exhaustive and theoretically grounded sets of security and privacy concerns, misconceptions, and principles, as well as broadening the range of stakeholders and contexts. 

Second, we plan to more systematically document and analyze the design process itself. In this work, ideation and refinement were conducted through iterative discussion and selection; future work will more explicitly capture this process through structured design tools (e.g., ideation cards or generative frameworks) and reflective documentation. This will enable us to better articulate how concepts are generated, filtered, and developed, addressing current limitations in methodological transparency and supporting adoption by other researchers. For example, we have begin to develop structured ideation cards representing concepts, misconceptions, domains, and interaction modalities. These cards will be combined in small sets to scaffold exploration and examine how constrained prompts shape design diversity, thematic coverage, and narrative plausibility. We will also explore generative AI as an ideation partner (the designs in this paper were developed without AI assistance). To better understand these processes, we plan to document the design work through autoethnography to identify recurring practices and challenges, including topic selection and balancing plausibility with provocation. 

Third, we aim to evaluate how these vignettes function in practice in user studies. While this paper positions them as communicative probes, our future work will examine how different audiences interpret and engage with the designs, such software develpers and other non-specialists. In particular, we are interested in how vignettes shape users’ threat identification of security and privacy risks, what kinds of discussions they provoke, and whether they reveal gaps, misconceptions, or alternative perspectives that are not easily captured through traditional methods.

Finally, we will explore how critical design vignettes can be integrated into existing security practices, including threat modelling and design ideation workflows. Given the inherently speculative nature of threat identification, we see potential for these artefacts to complement conventional approaches by supporting early-stage exploration of ``what can go wrong'', as suggested by Shostack + Associates \cite{shostack_fast_cheap_good}, in ways that are accessible to non-specialists. Future work will investigate how such methods can be incorporated into interdisciplinary teams and whether they can meaningfully inform design and decision-making processes. Through these directions, we aim to move from an initial corpus of designs toward a more robust methodological framework for using speculative artefacts as tools for reflection, communication, and inquiry within usable security and HCI research.





.


%\section{Concluding Remarks}

%\textblue{To be completed}





% = = = = = Ack = = = = = %

\begin{anonsuppress}
\begin{acks}
%We thank the reviewers who helped to improve our paper. 
%
J. Clark acknowledges support for this research project from the \grantsponsor{nserc}{National Sciences and Engineering Research Council of Canada (NSERC)}{https://www.nserc-crsng.gc.ca} through the NSERC, Raymond Chabot Grant Thornton, and Catallaxy Industrial Research Chair in Blockchain Technologies \grantnum[https://www.nserc-crsng.gc.ca/Chairholders-TitulairesDeChaire/Chairholder-Titulaire_eng.asp?pid=1045]{nserc}{ IRCPJ/545498-2018} and a Discovery Grant \grantnum{nserc}{RGPIN/04019-2021}
%
%We also acknowledge partial support from 
%\grantsponsor{amf}{Autorité des Marchés Financiers (AMF)}{https://lautorite.qc.ca} and
%\grantsponsor{chaire}{The Chaire Fintech: AMF -- Finance Montréal}{https://chairefintech.uqam.ca} and
%\grantsponsor{opc}{Office of the Privacy Commissioner of Canada (OPC)}{https://www.priv.gc.ca}.
%
\end{acks}
\end{anonsuppress}

% = = = = = Bibliography = = = = = %
\bibliographystyle{ACM-Reference-Format}
\bibliography{bib/design.bib}

% = = = = = End Notes = = = = = %

\end{document}

